\scanpage{58}
\mspage{45}
% archive sheet number: 57

dans $M(I^{\unsure{\circ|0}})$, considérons les \unsure{triples|4-tuples|k-tuples}
\[
 (J,Z_{\deleted{\unsure{J}}},\deleted{\unsure{p}}i_{\deleted{\unsure{J}}},p_{\deleted{\unsure{J}}})
\]

où $J\subset I$ est un ouvert, $Z_{\deleted{\unsure{J}}}$ un objet de $M(J^\circ)$, et $i_{\deleted{\unsure{J}}},p_{\deleted{\unsure{J}}}$ sont \unsure{des|les} flèches
\[
 \left\{\begin{aligned}
 &X|J\xrightarrow{i} Z\xrightarrow{p}Y|{\unsure{J}}\\
 &\text{avec}\quad i\in C_{J^\circ},\quad p\in W_{J^\circ}.
 \end{aligned}\right.
\]

On ordonne les quadruples par prolongement. Si on a un ens. \unsure{tot.} ordonné, de $(J_\alpha,Z_\alpha,\deleted{\unsure{f}}i_\alpha,p_\alpha)$, posant $J=\bigcup J_\alpha$, on trouve des $Z,i,p$ \deleted{\unsure{qui}} \inserted{uniques} \unsure{sur} $J$ qui prolongent les $Z_\alpha,i_\alpha,p_\alpha$ donnés, de plus on aura $i\in C_{J^\circ},p\in W_{J^\circ}$, car l'appartenance à \deleted{\unsure{des classes}} $C_{J^\circ}$ resp. $W_{J^\circ}$ est question locale sur $J$. Donc \inserted{par Zorn,} on est ramené à prouver qu'\deleted{il y a un} si $(J,Z,i,p)$ est maximal, on a $J=I$ -- d'où la factorisation voulue $f=pi$.

Donc : prouver que si $J\ne I$, le quadruple n'est pas maximal. Considérons le fermé \inserted{non vide} $I\setminus J$ \deleted{$\unsure{=A}$}. $I$ est supposé artinien, il y a un \uline{élément minimal} $\unsure{i_0}$ de $I\setminus J$. Posons
